\section{The Architect's Question}\label{the-architects-question}

After this chapter we should be able to reason about what a model's
parameter structure means for deployment: dense versus
Mixture-of-Experts (MoE), total parameters versus active parameters per
token, and what each implies for compute budgets and memory budgets. We
will meet the mechanisms (how MoE routes, how attention weights are
computed) not to learn how to train a model but so that, as architects,
we can look at a model card and immediately know whether the quoted
parameter count is what we actually pay for at inference time, and
whether the KV cache will grow as expected. After this chapter we can
ask: \emph{given a workload and a model architecture, what is the real
residency and compute cost per token?}

\section{Concept}\label{concept}

The \textbf{parameter allocation regime} of a model determines how many
of its weights are touched on each forward pass. This regime is not a
property of the model family name alone; it is a consequence of the
training objective and the routing mechanism.

\begin{itemize}
\tightlist
\item
  \textbf{Dense models} activate \emph{all} parameters for every token.
  A 70B dense model reads all 70B weights (or close to it) for every
  token, because every layer's full hidden dimension participates in
  every attention and feed-forward computation. The parameter count on
  the model card is the inference cost count.
\item
  \textbf{Mixture-of-Experts (MoE) models} route each token to a
  \emph{fraction} of the total parameters. A MoE model may have 236B
  total parameters distributed across eight experts, but only the top-2
  (or top-1) experts are activated per token. The quoted total parameter
  count includes the unused experts; the \emph{active} parameter count
  per token is the fraction that routing selects.
\end{itemize}

The architect-relevant distinction: \textbf{total parameters ≠ active
parameters per token.} This inequality is the single most important
number to read from a model card when sizing a deployment, because it
directly changes the compute cost per token. However, the KV cache does
not share this privilege --- we will show why MoE sparsity saves compute
but not memory.

A note on provenance: the total-versus-active distinction arises from
the training regime. In dense pre-training, every token sees the full
model gradient; in MoE pre-training, a routing loss directs each token
to the most relevant expert(s), and the unused experts receive no
gradient for that step. The routing mechanism itself (top-k, importance
sampling, learned gates) is an architectural decision that outlives
training and becomes a deployment fact.

\section{Mental Model}\label{mental-model}

Think of a model's parameters as a \textbf{resource pool} that is
allocated differently depending on the routing regime.

\begin{itemize}
\tightlist
\item
  In a \textbf{dense} model, the pool is \emph{fully allocated} every
  time. If the model is 70B params in FP16, we need 140 GB of weight
  memory available for every token. There is no ``spending less'' ---
  the full 140 GB is the residency floor.
\item
  In an \textbf{MoE} model, the pool is \emph{fractionally allocated}.
  The total pool may be large (reflecting the model's capacity and
  training compute), but what we actually wire up per token is a slice.
  With top-2 routing from 8 experts, roughly 2/8 = 25\% of the total
  parameter mass is active per token. This fraction saves compute (fewer
  FLOPs) but the attention layers still see every token, so the KV cache
  still grows with context length just as in a dense model.
\end{itemize}

The durable mental model: \textbf{total parameters set the ceiling;
active parameters per token set the floor.} The architect must track
both numbers.

\section{Worked Example}\label{worked-example}

We now carry concrete arithmetic using the canonical enterprise Q\&A
scenario (§14: \textasciitilde2,000 users, \textasciitilde10 rps
average, \textasciitilde40 rps peak, 1,200 + 8K context
\textasciitilde9.2K input, 300-token output, 70B-class model, FP16, 1
host 8×H100). The numbers below are worked out explicitly; where a
model's exact spec is not first-party verified, it is labeled
{[}VERIFY{]}.

\subsection{Dense 70B model (canonical,
{[}1P{]})}\label{dense-70b-model-canonical-1p}

\begin{longtable}[]{@{}
  >{\raggedright\arraybackslash}p{0.3333\linewidth}
  >{\raggedright\arraybackslash}p{0.3333\linewidth}
  >{\raggedright\arraybackslash}p{0.3333\linewidth}@{}}
\toprule\noalign{}
\begin{minipage}[b]{\linewidth}\raggedright
Metric
\end{minipage} & \begin{minipage}[b]{\linewidth}\raggedright
Value
\end{minipage} & \begin{minipage}[b]{\linewidth}\raggedright
Derivation
\end{minipage} \\
\midrule\noalign{}
\endhead
\bottomrule\noalign{}
\endlastfoot
Total parameters & 70B & §14 canonical scenario {[}1P{]} \\
FP16 residency (total) & 140 GB & 70B × 2 bytes = 140 GB {[}1P
DERIVED{]} \\
Active parameters per token & 70B & All parameters active for every
token {[}1P FACT{]} \\
FLOPs per forward pass (approx.) & \textasciitilde2.8T & 70B × 4
FLOP/param typical for transformer {[}2° DERIVED{]} \\

\label{tab:3.1}\end{longtable}

\emph{Table 3.1 --- Dense 70B model parameter arithmetic (worked
example, not reference).}

\subsection{MoE model representative (Mixtral 8x7B class,
{[}2°{]})}\label{moe-model-representative-mixtral-8x7b-class-2}

We contrast with a representative MoE model in the Mixtral lineage. The
exact parameter split is not always published in vendor model cards; the
numbers below are plausibility-checked and flagged.

\begin{longtable}[]{@{}
  >{\raggedright\arraybackslash}p{0.3333\linewidth}
  >{\raggedright\arraybackslash}p{0.3333\linewidth}
  >{\raggedright\arraybackslash}p{0.3333\linewidth}@{}}
\toprule\noalign{}
\begin{minipage}[b]{\linewidth}\raggedright
Metric
\end{minipage} & \begin{minipage}[b]{\linewidth}\raggedright
Value
\end{minipage} & \begin{minipage}[b]{\linewidth}\raggedright
Derivation
\end{minipage} \\
\midrule\noalign{}
\endhead
\bottomrule\noalign{}
\endlastfoot
Total parameters (across 8 experts) & \textasciitilde47B &
\textasciitilde6B per expert × 8 + embeddings {[}2° DERIVED{]} \\
Active parameters per token (top-2 routing) & \textasciitilde14B & Top-2
from 8 experts, \textasciitilde3--5\% active/total ratio (frontier MoE
{[}2° DERIVED{]}); actual activation \textasciitilde25\% for
top-2-from-8 configuration shown \\
FP16 active residency per token & \textasciitilde28 GB & 14B × 2 bytes =
28 GB {[}VERIFY DERIVED{]} \\
FLOPs per forward pass (approx.) & \textasciitilde0.56T & 14B active × 4
FLOP/param {[}VERIFY DERIVED{]} \\
Compute ratio (active/total) & \textasciitilde3--5\% (frontier) & 14B ÷
47B {[}2° DERIVED{]} \\

\label{tab:3.2}\end{longtable}

\emph{Table 3.2 --- MoE model parameter arithmetic (worked example, not
reference).}

\subsection{Why MoE saves compute but not KV-cache
memory}\label{why-moe-saves-compute-but-not-kv-cache-memory}

The compute savings are straightforward: if only \(n_a = 14\)B of
\(n_t = 47\)B total parameters are active per token, the FLOP count
drops by the active fraction. Since per-token FLOPs scale with the
number of parameters touched,

\[
\text{FLOP}_{\text{MoE/token}} = \frac{n_a}{n_t} \times \text{FLOP}_{\text{dense/token}} = \frac{14}{47} \times 2.8 \text{ T} \approx 0.56\text{ T}
\]

--- roughly a 5× reduction in compute per token. This is why MoE models
can be ``larger'' (more total parameters for the same training FLOP
budget) while keeping per-token cost comparable to a smaller dense
model.

The KV-cache story is the surprising part. Recall from \hyperref[chap:1]{Chapter 1} (§KV
cache as a concept, \hyperref[chap:1]{Chapter 1}) that the KV cache stores one key and one
value tensor per token, per layer. The cache size per token scales with
the canonical \hyperref[chap:1]{Chapter 1} formula,
\(KV_{\text{per-token}} = 2 \times n_\text{layers} \times d_\text{hidden} \times \text{bytes-per-value}\).
Critically, \textbf{the attention mechanism processes every token in the
context densely} --- each token's query attends to all previous tokens'
keys and values, regardless of whether the model is dense or MoE. MoE
sparsity operates in the feed-forward sub-layer; it does not change the
attention sub-layer's behavior.

Therefore, for a 70B-class model (dense or MoE) with 8K--9.2K input
context:

\begin{itemize}
\tightlist
\item
  \textbf{Dense 70B}: KV cache \textasciitilde1.3 MB per token × 9,200
  tokens ≈ 12 GB (at 8-bit) or \textasciitilde24 GB (at FP16-relevant
  precision for the key/value storage pattern discussed in Ch. 7).
\item
  \textbf{MoE 70B-equivalent}: The KV cache is \emph{identical} in size,
  because attention still processes all 9,200 tokens densely. The fact
  that only 14B of 47B total parameters are active in the feed-forward
  layers does not reduce the number of keys/values emitted by the
  attention layers.
\end{itemize}

In other words: MoE sparsity = compute sparsity (fewer FLOPs per token).
MoE sparsity ≠ memory sparsity (KV cache still grows linearly with
context length, same as dense). Only changing the attention mechanism
itself (e.g., hybrid/linear attention, compressed latent attention, or
KV offload) can stop the cache from growing. This distinction ---
compute-sparsity versus memory-sparsity --- is one an architect has to
get exactly right.

\begin{quote}
\textbf{Key takeaway:} MoE lets us fit a larger total parameter budget
within a weight-residency constraint (active weights take less FP16
memory), but the KV-cache budget must be provisioned as if the model
were dense. Do not assume MoE reduces KV-cache memory.
\end{quote}

\section{Measurement}\label{measurement}

How do we determine the total-versus-active split for a model we haven't
trained?

\begin{enumerate}
\def\labelenumi{\arabic{enumi}.}
\item
  \textbf{Model card inspection.} Most model vendors publish the total
  parameter count. The active-per-token count is harder to find; it
  requires reading the architecture documentation for the routing
  mechanism (top-k, top-1, importance-based). If the model card says
  ``70B parameters'' without qualification, assume dense unless it
  explicitly describes an MoE layout.
\item
  \textbf{Runtime measurement.} For a given input, profile the number of
  FLOPs or the memory footprint of active weights. Tools such as
  \texttt{torch.profiler} can count active parameters, but this is an
  empirical measurement, not a first-party spec.
\item
  \textbf{Architecture diagram.} For MoE models, the number of experts ×
  parameters-per-expert gives the total; the routing hyperparameter
  (top-k) × parameters-per-expert gives the active per token. This is a
  derived fact from the source code/release, not always on the model
  card.
\end{enumerate}

\begin{quote}
\textbf{Practical habit:} When we encounter a model quoted in ``X
billion parameters,'' pause and ask: \emph{is this total or active?} For
dense models they are the same; for MoE models they diverge. Misreading
this is the most common source of KV-cache and compute underestimation.
\end{quote}

\section{Common Mistakes}\label{common-mistakes}

\begin{itemize}
\item
  \textbf{Assuming total parameters = active parameters per token.} This
  is the single most costly misread. For a MoE model, the quoted
  ``236B'' or ``47B'' is the total across all experts; the active count
  per token is a fraction. Using the total number to size KV cache or
  memory will overestimate compute needs (if we think we need 236B × 2
  bytes we'll over-provision weight memory) or underestimate it (if we
  size for active but forget KV cache is dense).
\item
  \textbf{Confusing compute sparsity with memory sparsity.} MoE reduces
  FLOPs per token because fewer parameters are active. It does
  \emph{not} reduce the number of keys/values written to the KV cache.
  An architect who assumes MoE shrinks KV cache will be blindsided by
  memory pressure at long context lengths.
\item
  \textbf{Using MoE as a free lunch for long-context workloads.} The
  compute savings from MoE are real, but the memory cost (weights + KV
  cache) must be provisioned at the total-parameter scale for weight
  residency and at the-full-context scale for KV cache. Neither benefit
  is ``free'' in the other domain.
\item
  \textbf{Ignoring routing overhead.} Top-k routing itself requires
  computing routing scores (a softmax over expert assignments). This
  adds a small but non-zero compute cost that is sometimes omitted from
  published FLOP counts. It is usually negligible compared to the
  feed-forward arithmetic but should be acknowledged.
\end{itemize}

\section{Architecture Consequence}\label{architecture-consequence}

Knowing whether a model is dense or MoE, and whether the quoted
parameter count is total or active, directly changes three architectural
decisions:

\begin{enumerate}
\def\labelenumi{\arabic{enumi}.}
\item
  \textbf{Weight residency budget.} For a dense 70B model, provision 140
  GB of FP16 weight memory per host. For a MoE model with
  \textasciitilde47B total and \textasciitilde14B active, provision
  \textasciitilde28 GB for the active weights + a smaller overhead for
  the unused experts (they may be swapped or kept on device depending on
  the deployment framework). MoE allows a larger total parameter budget
  within the same weight-memory ceiling, but the active-fraction budget
  is what matters for ``weights on device.''
\item
  \textbf{KV-cache provisioning.} Provision the KV cache as for a dense
  model of the same context length and hidden dimension. Do not apply a
  MoE sparsity factor to the KV cache. If we are running a 9.2K-token
  input on a 70B-class model (dense or MoE), the KV cache memory is the
  same order of magnitude --- plan accordingly.
\item
  \textbf{Parallelism and routing strategy.} MoE deployment requires a
  routing strategy (e.g., expert placement, dispatcher, switchboard) and
  often expert parallelism across GPUs. This adds infrastructure
  complexity (cross-GPU communication for expert routing) that dense
  models do not have. The architectural choice between dense and MoE is
  therefore not just a parameter-count decision; it is a
  deployment-complexity decision.
\end{enumerate}

\subsection{The 2026 baseline shift: same framework, new
constants}\label{the-2026-baseline-shift-same-framework-new-constants}

\textbf{The 2026 trend signal (forward-looking).} By August 2026 the
four frontier open-weight families converge on one direction, and it is
not more of the same dense scaling: extreme MoE sparsity (single-digit
active-parameter fractions --- Qwen 6 B of 125 B, Kimi K3 104 B of 2.8
T) coupled with \textbf{hybrid attention} (compressed/sparse/linear
hybrids such as Gated DeltaNet + sparse attention) has become the
standard way to make long-context inference affordable. \hyperref[chap:27]{Chapter 27}
(Appendix A) documents these at full depth with first-party provenance;
what matters here, in the model-understanding chapter, is that this is
not a competing mechanism but the \emph{same} dense-vs-MoE,
total-vs-active, KV-vs-compute machinery this chapter just taught ---
with new constants. We hold that view explicitly as the
\textbf{forward-looking direction} the architect should size toward, and
flag it below as continuing future work (the frontier moves; this
handbook is a living document, and Appendix A is where the moving part
lives).

\textbf{Why the trend signal should not be read as ``dense is dead''.}
The convergence of the four \emph{frontier} families on extreme MoE
sparsity describes the top of the market --- the
\textasciitilde100B-and-up tier that anchor
\textasciitilde2.8T-parameter fleets. It is not a claim that dense
models have disappeared. On the contrary, the smaller, denser tier is
alive and actively shipping in 2026 for exactly the
constrained-deployment reasons this chapter teaches:
\textbf{Qwen3.8-27B} ({[}1P: HF Qwen/Qwen3.8-27B{]}) is a dense 27.8B
model whose 4-bit weights fit on a single 24 GB consumer GPU, and
\textbf{Muse Glimmer 30B} ({[}1P: Meta via vLLM-Recipes{]}) is a dense
29.6B local-agentic model that runs in \textasciitilde18 GB. These exist
because the dense-vs-MoE trade is itself a \emph{deployment} decision,
not a date: MoE buys active-parameter efficiency at the cost of
total-parameter memory and routing complexity --- worthwhile when you
serve many requests from a large pool of resident weights, but hard to
justify when your whole budget is one consumer GPU whose total weights
already overrun. That is precisely why this book keeps its canonical on
the dense tier: the canonical must serve the architect whose \emph{whole
fleet} is a handful of boxes, not only the one sizing a 2.8T datacenter
model. The framework does not change between tiers --- total vs active,
KV per-token, and residency budgeting apply identically to a 27B dense
laptop model and a 125B/6B MoE fleet model.

\textbf{Why the rest of this chapter's canonical stays a dense full-MHA
model.} If the trend is MoE + hybrid, why does this book's canonical
still size a dense 70B full-MHA workload? Because a dense model with
full multi-head attention gives the cleanest single-constant KV formula
(\texttt{2\ ×\ layers\ ×\ hidden\ ×\ bytes}) and the largest KV
footprint --- the conservative worst case an architect should always
size against first. Teaching against that bound means the reader learns
the framework on the hardest memory case and then relaxes it. It is a
deliberate teaching simplification, not a claim that in 2026 an
enterprise RAG fleet really runs on dense 70B.

\textbf{One verified 2026 transfer: Qwen3.8-Flash-Next.} The model card
{[}1P: HF Qwen/Qwen3.8-Flash-Next{]} is the Qwen4-architecture preview:
\textbf{125 B total / 6 B active} (plus a separate 51 B of offloadable
N-gram embedding parameters), \textbf{512 MoE experts} (10 routed + 1
shared per token), and \textbf{hybrid attention} --- Gated DeltaNet
(GDN) on three of every four layers plus Qwen Sparse Attention (QSA) on
the fourth. This is exactly the ``125B/6B active MoE + Gated
DeltaNet/QSA'' class of model a 2026 architect actually considers.
Stepping through this book's own arithmetic:

\begin{itemize}
\tightlist
\item
  \textbf{Weight residency} (Ch. 7's floor): active weights are
  \textasciitilde6 B × 2 B = \textasciitilde12 GB, not the canonical 140
  GB --- MoE sparsity collapsed the on-GPU weight floor by
  \textasciitilde12×. But total weights dominate the memory ceiling:
  \textasciitilde125 B × 2 B ≈ 250 GB (plus the offloadable N-gram), so
  the \emph{total-vs-active} split of this chapter is not a footnote ---
  it is the difference between fitting on one host and sharding across
  several.
\item
  \textbf{KV cache} (Ch. 7/8): the canonical
  \texttt{2\ ×\ layers\ ×\ hidden\ ×\ bytes} KV formula is the full-MHA
  bound. Qwen's hybrid attention attacks the KV constant itself at the
  mechanism level --- GDN keeps a fixed-size recurrent state (context
  cost no longer grows linearly per token across those layers) and QSA
  attends sparsely at micro-block granularity --- so per-token KV no
  longer follows the dense formula. The framework question (``does it
  fit?'') is unchanged; the \emph{function} you plug in for KV is not
  the dense one.
\item
  \textbf{Compute} (Ch. 8): active-parameter FLOPs drop to
  \textasciitilde2 × 6 B per token, but the QSA/GDN layers ride
  different roofline points than a dense 70B prefill/decode split, so
  the throttled-resource conclusion must be re-derived per model (Ch.
  8's measurement discipline, not a one-time number).
\end{itemize}

\textbf{The architect's move is unchanged.} For the canonical dense
model this book sizes memory, compute, and fleet with one clean constant
set. For a 2026 MoE/hybrid model the architect does the \emph{same}
decision loop --- read the model card, get total vs active, get the KV
scheme's per-token function, size residency and roofline, then decide
parallelism (Ch. 10) and serving (Ch. 11). Only the constants differ;
the framework does not. \hyperref[chap:27]{Chapter 27} walks the four frontier models
through exactly this re-derivation and flags which vendor numbers still
carry a verify-on-own-hardware caveat.

\begin{quote}
\textbf{Recorded as future work.} This forward-looking view is
intentionally a \emph{pointer}, not the full treatment --- the complete,
current analysis of the 2026 frontier lives in \hyperref[chap:27]{Chapter 27} (Appendix A).
Because the frontier moves quickly and every vendor figure in it carries
a verify-on-own-hardware caveat, we treat the trend as a living section
to be updated against new model cards, rather than a set of frozen
gospel numbers (per the handbook's ``living document'' framing,
book-architecture.md §10). An architect should re-check Appendix A
before any capacity decision --- the constants, not the framework, are
what change.
\end{quote}

\section{What We Still Don't Know}\label{what-we-still-dont-know}

As of 2026-08, several questions remain open and are flagged
{[}VERIFY{]}:

\begin{itemize}
\tightlist
\item
  \textbf{Exact active-fraction per token for commercial MoE models.}
  Model cards rarely publish the precise top-k routing distribution
  under real workloads. We know the architectural top-k (e.g., top-2
  from 8 experts), but the actual fraction of active parameters may vary
  with prompt content, token position, and expert availability.
  {[}VERIFY HYPOTHESIS{]}
\item
  \textbf{Whether router latency offsets MoE compute savings.} The cost
  of computing routing scores and dispatching to experts is not always
  included in published FLOP counts. On some hardware the routing
  overhead can be significant enough to narrow the compute gap between
  MoE and dense. {[}VERIFY HYPOTHESIS{]}
\item
  \textbf{KV-cache behavior under MoE with expert swapping.} When
  experts do not fit on a single GPU and must be swapped (offloaded),
  does the KV cache interact with the swapping mechanism in ways that
  change its effective size or access pattern? This has not been
  systematically characterized. {[}VERIFY HYPOTHESIS{]}
\end{itemize}

These flags exist because home-lab measurements are not reference per
the evidence taxonomy; they will be resolved (promoted, dropped, or
demoted) before publication.

\section{End-of-Chapter Mini-Case}\label{end-of-chapter-mini-case}

\emph{(Continuous scenario --- this is its first appearance, chaining
from the enterprise Q\&A thread.)}

An architect is fleshing out the design of the internal Q\&A tool
described in \hyperref[chap:1]{Chapter 1}'s mini-case. The team is torn between a dense 70B
model and a MoE model with 8 experts (adisedly total \textasciitilde47B
params, \textasciitilde14B active per token). They ask: ``If we go MoE,
can we cut our GPU memory budget?''

From the model-understanding chapter, the architect can answer:
\emph{Weight residency yes --- active parameters take \textasciitilde28
GB vs.~140 GB for a dense 70B, we could fit the model on fewer GPUs or a
smaller host. But KV-cache memory no --- the attention layers still
process every token densely, so the cache budget for a 9.2K-context
input is the same regardless of whether the model is dense or MoE. We
still need to provision for \textasciitilde12 GB of KV cache (at 8-bit)
or \textasciitilde24 GB (at the precision pattern discussed in Ch. 7)
per request, and we'll also need expert-routing infrastructure on top of
that.}

The architect's decision therefore hinges on whether the compute savings
from MoE (fewer FLOPs per token, potentially lower
\(/token) outweigh the added routing complexity and the fact that KV-cache memory is unchanged. The team decides on a MoE model for the compute/\)
advantage, but provisions the KV-cache budget at the dense-model rate,
and adds one GPU dedicated to the routing dispatcher.

\begin{center}\rule{0.5\linewidth}{0.5pt}\end{center}

\begin{figure}
\centering
\pandocbounded{\includegraphics[keepaspectratio,alt={\hyperref[fig:3.1]{Fig 3.1} --- Dense vs MoE parameter allocation and KV-cache behavior {[}2° DERIVED{]}},width=\maxwidth{\textwidth},keepaspectratio]{figures/fig-03-0301.pdf}}
\caption{\hyperref[fig:3.1]{Fig 3.1} --- Dense vs MoE parameter allocation and KV-cache
behavior {[}2° DERIVED{]}}

\label{fig:3.1}\end{figure}

\emph{\hyperref[fig:3.1]{Fig 3.1} --- Parameter-activation contrast. A dense 70B activates
all 70B parameters per token (residency ≈ 140 GB, KV grows linearly with
context). An 8-expert MoE activates only the top-2 (\textasciitilde14B)
per token (residency ≈ 28 GB active), but KV-cache growth with context
is identical to dense --- attention still processes every token.}

\subsection{\hyperref[tab:3.1]{Table 3-1} --- Parameter and KV-cache arithmetic for the
canonical
workload}\label{table-3-1-parameter-and-kv-cache-arithmetic-for-the-canonical-workload}

\begin{longtable}[]{@{}
  >{\raggedright\arraybackslash}p{0.2500\linewidth}
  >{\raggedright\arraybackslash}p{0.2500\linewidth}
  >{\raggedright\arraybackslash}p{0.2500\linewidth}
  >{\raggedright\arraybackslash}p{0.2500\linewidth}@{}}
\toprule\noalign{}
\begin{minipage}[b]{\linewidth}\raggedright
metric
\end{minipage} & \begin{minipage}[b]{\linewidth}\raggedright
dense 70B
\end{minipage} & \begin{minipage}[b]{\linewidth}\raggedright
MoE (8×7B class)
\end{minipage} & \begin{minipage}[b]{\linewidth}\raggedright
derivation
\end{minipage} \\
\midrule\noalign{}
\endhead
\bottomrule\noalign{}
\endlastfoot
total parameters & 70B & \textasciitilde47B & {[}2°{]} expert split \\
active params per token & 70B & \textasciitilde14B & top-2 from 8
{[}\textasciitilde3--5\% frontier MoE {[}2° DERIVED{]}; Mistral
top-2-from-8 \textasciitilde25\% for reference{]} \\
FP16 residency (total) & 140 GB & \textasciitilde94 GB & 70B × 2 / 47B ×
2 \\
FP16 residency (active) & 140 GB & \textasciitilde28 GB & 70B × 2 / 14B
× 2 \\
KV cache per request (9.2K ctx, 8-bit) & \textasciitilde12 GB &
\textasciitilde12 GB & identical --- attention dense \\
FLOPs per forward pass & \textasciitilde2.8T & \textasciitilde0.56T &
70B × 4 / 14B × 4 \\
compute speedup (active/total) & 1× & \textasciitilde5× & 14B ÷ 70B
{[}\textasciitilde3--5\% frontier; Mistral top-2-from-8
\textasciitilde25\% for reference{]} \\

\label{tab:3.3}\end{longtable}

\emph{All figures trace to the §14 canonical scenario; MoE column is a
{[}2°{]} worked variant, not a measurement claim.}

\begin{center}\rule{0.5\linewidth}{0.5pt}\end{center}

\emph{Next chapter: \hyperref[chap:4]{Chapter 4} --- The Anatomy of an AI Workload, which
characterizes the enterprise Q\&A RAG workload across the six dimensions
(quality, traffic, token profile, latency, economic, operational
constraints) and uses the same continuous mini-case scenario.}

\emph{Canonical scenario citation: all numerical values in this chapter
that trace to the §14 canonical enterprise Q\&A RAG workload are drawn
from the fixed source of truth defined in book-architecture.md §14. The
dense 70B numbers are {[}1P{]} per the canonical scenario; MoE numbers
are {[}2°{]} plausibility-checked variants.}
